%

\section{Aksjomaty Tarskiego}
\index{aksjomat!Tarskiego}%

Zmienne są punktami, relacja $\beta$ mówi, że punkty leżą na prostej w kolejności takiej, jak w napisie, zaś relacja $\equiv$ mówi, że dwa odcinki są przystające:
\begin{enumerate}
    \item jeśli $\beta (x, y, x)$, to $x = y$;
    \item jeśli $\beta (x, y, u)$ oraz $\beta (y, z, u)$,  to $\beta (x, z, u)$;
    \item jeśli $\beta (x, y, z)$, $\beta (x, y, u)$ oraz $x \neq y$, to $\beta (x, z, u)$ lub $\beta (x, u, z)$;
    \item $xy \equiv yx$;
    \item jeśli $xy \equiv zu$ oraz $xy \equiv vw$, to $zu \equiv vw$;
    \item jeśli $xy \equiv zz$, to $x = y$;
    \item aksjomat Pascha: jeśli $\beta (x, t, u)$, $\beta (y, u, z)$, to istnieje taki $v$, że $\beta (x, v, y)$, $\beta(z, t, v)$'';
    \item V aksjomat Euklidesa\footnote{Ten aksjomat można zastąpić przez ,,na każdym trójkącie można opisać okrąg'' (mniej zmiennych) albo ,,linia środkowa w trójkącie jest dwukrotnie krótsza od podstawy'' (brak kwantyfikatorów).} (dla każdego kąta i punktu wewnątrz istnieje odcinek przez ten punkt, który przecina ramiona kąta): jeśli $\beta(x, u, t)$, $\beta (y, u, z)$, $x \neq y$, to istnieją takie $v, w$, że $\beta(x, z, v)$, $\beta(x, y, w)$ oraz $\beta (v, t, w)$;
    \item aksjomat pięciu odcinków (cecha przystawania bkb): jeśli $xy \equiv x'y'$, $yz \equiv y'z'$, $xu \equiv x'u'$, $yu \equiv y'u'$, $\beta (x, y, z)$, $\beta(x', y', z')$ i $x \neq y$, to $zu \equiv z'u'$;
    \item odkładanie odcinków: istnieje $z$ taki, że $\beta (x, y, z)$ i $yz \equiv uw$;
    \item istnieją trzy niewspółliniowe punkty\footnote{Bez tego aksjomatu teoria ma model na prostej rzeczywistej.} $x, y, z$: $\neg \beta(x, y, z)$, $\neg \beta(y, z, x)$, $\neg \beta(z, x, y, )$;
    \item trzy punkty równoodległe od końców odcinka tworzą prostą\footnote{Bez tego aksjomatu teoria ma model w trzech i więcej wymiarach.}: jeśli $xu \equiv xv$, $yu \equiv yv$, $zu \equiv zv$ i $u \neq v$, to $\beta(x, y, z)$, $\beta(y, z, x)$ lub $\beta(z, x, y)$;
    \item schemat aksjomatu ciągłości.
\end{enumerate}

% TODO: https://en.wikipedia.org/wiki/Tarski%27s_axioms

% TODO: https://smp.uws.edu.pl/msn/14/10-18.pdf
% Kompletne, beznadziejne bezdroża, Marek KORDOS

% https://www.deltami.edu.pl/media/articles/2015/01/delta-2015-01-prozny-trud.pdf

Patrz też Greenberg \cite{greenberg_2010}. % TODO: czemu patrz?

%